\chapter{UML and Detailed Software Design}\label{ch:uml}

This chapter presents five UML views of the system. Each diagram was drawn from the source code named in the
comment at the top of its TikZ file under \code{report/figures/uml/}, and the class, operation and message names
are those used in the code. The diagrams are rendered by TikZ when the report is compiled, so the editable source
and the figure are the same file.

\section{Use Case Diagram}

\autoref{fig:uml-usecase} shows the system boundary, the primary actor (the researcher) and the two external
systems that take part in use cases: arXiv supplies metadata and PDF files, and the local Ollama service generates
answers and analyses. Importing or uploading a paper \emph{includes} tracking its ingestion job, because both
requests return a job that the interface polls. Retrying a failed job \emph{extends} job tracking, since it happens
only when a job has failed. Inspecting a citation \emph{extends} asking a question and viewing history: it is
optional behaviour available wherever an answer with citations is shown. The evaluation harness is a maintainer
command-line tool and is deliberately outside the boundary (\autoref{ch:requirements}).

\begin{figure}[htbp]
  \centering
  \fitfigure{figures/uml/use-case}
  \caption{UML use case diagram}\label{fig:uml-usecase}
\end{figure}

\section{Class Diagrams}

The static structure is split into two diagrams to keep them legible. \autoref{fig:uml-class-domain} shows the
persistent domain model from \code{backend/app/db/models.py}. A \code{Paper} owns at most one \code{Document} (the
PDF currently behind it), which owns its \code{Chunk}s; deleting a paper cascades to both (composition). An
\code{IngestionJob} produces at most one paper. Question history is modelled by \code{Query}, \code{Answer},
\code{AnswerEvidence} and \code{AnswerCitation}: each citation resolves to at most one evidence row of the same
answer (none if the label was invalid), and each evidence row is a snapshot of a chunk whose reference becomes
empty if the chunk is later deleted. \code{Analysis} stores the identifiers of the papers it covers rather than a
foreign key, shown as a dependency.

\begin{figure}[htbp]
  \centering
  \fitfigure{figures/uml/class-domain}
  \caption{UML class diagram, part 1: persistent domain model}\label{fig:uml-class-domain}
\end{figure}

\autoref{fig:uml-class-services} shows the service layer and the abstractions it depends on. The services depend on
interfaces (\code{GenerationProvider}, \code{Embedder}, \code{RerankerProtocol}, \code{ArxivSource}), realised by the
production classes; the tests substitute fakes for the same interfaces. Citation checking and the search queries
are module-level functions, shown as \guillemotleft utility\guillemotright{} classes. Only public operations are listed,
and \code{CorpusService}'s use of \code{FileStore} for deleting files is drawn; framework classes are omitted.

\begin{sidewaysfigure}
  \centering
  \fitwide{figures/uml/class-services}
  \caption{UML class diagram, part 2: services and abstractions}\label{fig:uml-class-services}
\end{sidewaysfigure}

\section{Sequence Diagram}

\autoref{fig:uml-sequence} follows a question from the interface to the answer and back, matching
\code{QAService.ask}. After validation the service retrieves candidates (embedding, dense and lexical SQL, rank
fusion and cross-encoder re-ranking inside \code{SearchService}), keeps passages whose cosine similarity is at least
0.30 and labels them. The outer \code{alt} fragment shows the early abstention without a model call. In the other
operand the model is called; the inner \code{alt} fragment separates provider failures (Ollama unreachable, a
timeout, or output that does not validate against the schema), which are persisted with status \code{error} and
returned as HTTP 503, 504 or 502, from the normal path, in which citations are checked and the outcome is persisted
as \code{answered} or \code{insufficient\_evidence}. The diagram ends with the user opening a citation, which the
interface resolves from the evidence already returned; no further request is made.

\begin{sidewaysfigure}
  \centering
  \fitwide{figures/uml/sequence}
  \caption{UML sequence diagram: grounded question answering}\label{fig:uml-sequence}
\end{sidewaysfigure}

\section{Activity Diagram}

\autoref{fig:uml-activity} models ingestion in two swimlanes. The request lane is synchronous: an arXiv identifier is
parsed, or an uploaded file is read with a size limit, checked (PDF signature, encryption, page limit) and hashed.
Invalid requests return 422 and duplicate uploads 409 before any job exists. Valid uploads are stored under their
hash, a job is created (or the in-flight job for the same source is returned) and handed to the job runner, and the
request returns 202 at the fork. The background lane follows \code{run\_job}: an arXiv job fetches metadata and
finishes immediately if that version is already indexed with the current chunker; otherwise the PDF is downloaded
(or the stored upload read), extracted page by page, chunked, embedded and written in one transaction. Any error
raised inside the interruptible region---for example a network failure, HTTP 429 from arXiv, a PDF without a text
layer, or an embedding failure---ends the job as \code{failed} with its message, and a rejected upload's file is
deleted.

\begin{figure}[htbp]
  \centering
  \fitfigure{figures/uml/activity}
  \caption{UML activity diagram: paper ingestion}\label{fig:uml-activity}
\end{figure}

\section{Communication Diagram}

\autoref{fig:uml-communication} shows the same interaction as \autoref{fig:uml-sequence} as a communication
(collaboration) diagram. It emphasises which objects are linked: the route talks only to \code{QAService}, which is
the single object that talks to the language model, the citation checker and the database for persistence, while
\code{SearchService} encapsulates the embedder, the database queries and the re-ranker. Message numbers are nested
to show that messages 1.1.1.1 to 1.1.1.3 happen during \code{search}, and guards mark the conditional messages:
generation happens only when evidence passes the threshold, and citation checking only when generation produced
schema-valid output.

\begin{figure}[htbp]
  \centering
  \fitfigure{figures/uml/communication}
  \caption{UML communication diagram: grounded question answering}\label{fig:uml-communication}
\end{figure}

\section{Consistency Between the Diagrams}

\autoref{tab:uml-consistency} cross-checks the dynamic diagrams against each other and against the code.

\begin{table}[htbp]
  \centering
  \caption{Correspondence of sequence and communication messages}\label{tab:uml-consistency}
  {\small
  \begin{tabularx}{\textwidth}{@{}llL@{}}
    \toprule
    \textbf{Sequence} & \textbf{Communication} & \textbf{Code} \\
    \midrule
    2, 3 & 1, 1.1 & \code{api/v1/qa.py: ask} $\rightarrow$ \code{QAService.ask} \\
    5 & 1.1.1 & \code{SearchService.search} \\
    6 & 1.1.1.1 & \code{SentenceTransformerEmbedder.embed} \\
    7 & 1.1.1.2 & \code{search\_hybrid\_chunks} (dense and BM25 SQL, RRF) \\
    8 & 1.1.1.3 & \code{CrossEncoderReranker.rerank} \\
    10b & 1.1.2 & \code{OllamaProvider.generate} \\
    11b & 1.1.3 & \code{check\_citations} \\
    10a, 11a, 12 & 1.1.4 & \code{QAService.\_persist} \\
    \bottomrule
  \end{tabularx}
  }
\end{table}
